
\section{Conclusion}
\label{sec:conclusion}

This paper described the design and bench-scale evaluation of \textbf{NirmalNode}, a
closed-loop Green IoT architecture integrating real-time microenvironmental sensing,
online incremental stream forecasting, predictive relay actuation, and localized
three-stage physical air purification. On a controlled synthetic benchmark ($N=499$
prequential samples), the blended SGD+PA ensemble (M8) achieved MAE~10.616~\ugsm,
$R^2=0.7334$, and $F_1=0.940$ for hazard-actuation, with a 68.2\% reduction from
cold-start tracking error through pure streaming adaptation. On a Bangladesh
public-dataset-derived evaluation stream, zero false negatives were observed in
$N=500$ samples ($F_1=0.992$, AUC-ROC~0.985). Laboratory bench trials recorded 95.7\%
single-pass PM$_{2.5}$ reduction and 98.6\% PM$_{10}$ reduction with a 4--6 second
pre-actuation observation; gas-phase results are qualitative MOS indicators only.
System peak power is 8.6~W with an indicative prototype component BOM of
USD~67.27 (BDT~7,400).

These results support the feasibility of the proposed integrated architecture for
localized occupational exposure monitoring and targeted air cleaning in
resource-constrained industrial settings. Field validation in genuine factory
environments, repeated filtration trials with reference instrumentation, predictive
vs.\ reactive controlled comparison, and multi-node deployment remain necessary
steps toward a deployable system.
